\begin{center}
\textbf{{\Large ABSTRACT}}\\
\begin{center}
\justifying
Log-based anomaly detection is a critical component of ensuring the reliability 
and security of large-scale software systems. Modern systems generate terabytes 
of log data daily, making automated detection of abnormal patterns essential for 
rapid incident response and failure prevention. Existing parsing-free methods such 
as NeuralLog have advanced the state of the art by eliminating dependency on brittle 
log parsers through BERT-based encoding of raw log text. However, these methods 
treat all tokens uniformly through simple averaging of contextual embeddings, 
discarding the structural organization of log messages, and remove numeric tokens 
that carry critical anomaly signals such as HTTP error codes and hardware identifiers. 
This project proposes a novel role-aware representation learning framework that 
addresses these limitations by inducing latent token roles — state, entity, action, 
quantity, and location — automatically from sequence-level anomaly labels, without 
any manual annotation or log parsing. A role induction network assigns a soft 
probability distribution over five role categories to each token, and role-aware 
pooling constructs structured message representations that preserve both semantic 
and structural information. Numeric tokens are explicitly retained to capture 
quantitative anomaly indicators. The framework is evaluated on four public benchmark 
datasets — HDFS, BGL, Thunderbird, and Spirit — achieving F1-scores of 0.98--0.99 
and consistently outperforming NeuralLog and all baseline methods. Ablation studies 
confirm that both role modeling and numeric token preservation contribute 
independently to the observed improvements. The learned token roles align with 
semantic expectations despite being induced solely from anomaly labels, providing 
an interpretable basis for anomaly attribution with negligible computational 
overhead over NeuralLog.

The second objective of this project investigates the vulnerability of XOR-based 
secret image sharing schemes to learning-based reconstruction attacks using 
autoencoders. Classical $n$-out-of-$n$ XOR schemes provide information-theoretic 
security against adversaries holding strict subsets of shares; however, their 
resilience against data-driven models that exploit learned image priors remains 
largely unexplored. This work proposes an encoder--decoder autoencoder framework 
that attempts to reconstruct a secret image from a constrained consecutive subset 
of its XOR-generated shares, without access to cryptographic keys or the remaining 
shares. The model leverages convolutional feature extraction and skip connections 
to recover structural content from the partial observations available to an attacker. 
Reconstruction fidelity is evaluated using PSNR and SSIM across multiple benchmark 
image datasets. Experimental results show that while meaningful reconstruction is 
possible when all shares are available, reconstruction fails under missing-share 
conditions, thereby preserving the theoretical security guarantees of the XOR scheme 
in practical settings.

The third objective of this project focuses on the cryptanalysis of a chaos-based 
medical image encryption scheme under the chosen-plaintext attack (CPA) model. 
The target scheme employs a permutation-diffusion architecture using Arnold's Cat 
Map and a Two-Dimensional Logistic Sine Coupling Map (2D-LSCM). Through structural 
analysis, this work demonstrates that the keystream used in the diffusion stage is 
independent of the plaintext, allowing it to be recovered exactly using a single 
chosen input. Furthermore, the separability of permutation and diffusion stages 
enables complete reconstruction of the permutation mapping using structured probe 
images. By combining these results, the entire encryption process can be inverted 
without knowledge of the secret key. Experimental validation confirms that the 
attack requires only $n+1$ oracle queries for an $n \times n$ image and achieves 
exact plaintext recovery with zero error. These findings reveal fundamental 
weaknesses in chaos-based encryption designs and highlight that statistical measures 
such as entropy and NPCR are insufficient indicators of cryptographic security.

\end{center}
\end{center}

\keywords{Log Anomaly Detection, Role-Aware Representation Learning, Token Role 
Induction, BERT, Parsing-Free Log Analysis, Deep Learning, Transformer, 
Representation Learning, Visual Cryptography, Secret Image Sharing, XOR Secret 
Sharing, Autoencoder, Reconstruction Attack, Adversarial Machine Learning, 
Image Security, Chaos-Based Encryption, Cryptanalysis, Chosen-Plaintext Attack}\\
\clearpage